\section*{Results}

\subsection*{Relational priors carry structure that flat priors discard}

Multi-omics resources released since 2020 --- tissue-specific eQTL
catalogues, protein--protein interaction networks at curated
confidence, base-pair-resolved regulatory-element atlases ---
encode far more relational structure than flat-prior fine-mappers
currently exploit. \textbf{Flat-prior fine-mapping collapses that
relational structure to a scalar.} In the SuSiE/FINEMAP/PolyFun family, biological annotations
enter as a per-variant weight vector $w \in \mathbb{R}^n$, averaged
across tissues and detached from the gene, pathway and interaction
context each variant belongs to. Constructing that scalar weight is
itself a substantial curation step: it requires harmonising
heterogeneous functional sources --- eQTL tables, regulatory-element
atlases, pathway curations and PPI networks --- into a single
per-variant weight, typically via stratified LD-score regression over
a pre-curated functional-category catalogue (the human Baseline 2.2
catalogue uses 96 such categories\cite{finucane2015partitioning}).
The burden is heaviest precisely where post-GWAS analysis is most
needed: in non-model species, including crops and livestock, where functional
annotations are sparse, fragmented and lack a pre-harmonised
catalogue. The graph representation never demands such a scalar ---
annotation layers register as typed edges where data exists, missing
edges stay missing, and end users get a ready-to-go pipeline without
species-specific harmonisation (see Discussion
``\nameref{disc:nonmodel}'' for the full argument). GraphGWAS instead represents this
context as a typed graph (Figure~\ref{fig:graph}): variants connect
to genes through $\texttt{HAS\_CONSEQUENCE}$ and tissue-specific
$\texttt{eQTL}$ edges, genes participate in pathways via
$\texttt{IN\_PATHWAY}$, and the gene layer is coupled within itself
by protein--protein $\texttt{INTERACTS\_WITH}$ edges. The graph loaded on the 1000 Genomes Phase~3 high-coverage cohort
\cite{1000genomes2015global,byrska2022high} carries 70.7~million
variants, 20{,}092 GENCODE~v47 genes\cite{frankish2023gencode},
43.2~million GTEx~v8 tissue eQTLs\cite{gtex2020}, 230{,}850 STRING
interactions at combined score~$\geq$~700\cite{szklarczyk2023string},
and 370{,}000 ENCODE cCRE regulatory elements\cite{encode2020expanded}
(full counts in Supplementary Table~S1). All
relational edges are first-class: a fine-mapping procedure can either
collapse them to a scalar prior (the flat-prior approach) or carry
them through the inference (the relational approach we develop
below).

\begin{figure}[htbp]
\centering
\includegraphics[width=0.95\textwidth]{fig1_architecture.pdf}
\caption{\textbf{Relational annotation priors carry biological
structure that flat priors collapse to a scalar.}
(\textbf{a})~Flat-prior fine-mappers (PolyFun / SuSiE family) accept
a vector of per-variant annotation weights $w_i$ that collapses
gene, tissue, pathway and interaction information into a single
scalar per variant; the Bayesian regression proceeds on
($z$-scores, LD matrix $R$, prior weights $w$) and returns a PIP
vector whose biological context must be reattached by a post-hoc
enrichment step.  (\textbf{b})~GraphGWAS exposes the multi-omics
knowledge graph directly: a query variant (highlighted red circle)
is connected by \emph{typed} edges---\texttt{HAS\_CONSEQUENCE} to
its coding genes, tissue-specific \texttt{eQTL} edges to the genes
it regulates, \texttt{IN\_REGULATORY} to ENCODE cCREs,
\texttt{IN\_PATHWAY} via genes to pathway nodes, and STRING
\texttt{INTERACTS\_WITH} edges coupling genes at combined
score~$\geq$~700---to five typed node classes (Variant, Gene,
Pathway, RegulatoryElement, Tissue; legend bottom-left). Real
labels are illustrative (FTO, IRX3, Adipose\_Subcutaneous, fatty-
acid metabolism). The credible-set output is itself a graph object:
each reported variant is co-queryable with its gene, tissue and
pathway neighbours in one traversal. The HBP / GAFM algorithm that
computes posterior inclusion probabilities by message passing
across this graph is detailed in Figure~\ref{fig:hbp_gafm}; this
panel shows the underlying data model, not the inference
procedure. Multi-omics graph state on 1000 Genomes Phase~3 (inline
lower-right box; full counts in Supplementary Table~S1): 70.7~M
variants, 20{,}092 GENCODE~v47 genes, 43.2~M GTEx~v8 eQTL edges
across 49 tissues, 230{,}850 STRING interactions, and 370{,}000
ENCODE cCRE regulatory elements.}
\label{fig:graph}
\end{figure}

\subsection*{Hierarchical belief propagation matches Bayesian baselines at 5--40$\times$ speed}

HBP converts the graph into iterated message passing across
variant, gene and pathway layers (Figure~\ref{fig:hbp_gafm}). Each
iteration runs an upward pass---variant beliefs aggregate to gene
scores weighted by eQTL strength, then diffuse across the gene layer
along STRING PPI edges at a 0.3 coupling, then aggregate to
pathway scores---followed by a downward pass that returns a variant
prior $\pi^{(t)}$ via the same edges. The variant belief is combined
as $b^{(t+1)} = \alpha \cdot \text{softmax}(u) + (1-\alpha) \pi^{(t)}$,
where $u$ is the LD-deconvolved z-statistic, and damped by
$\lambda$ before renormalisation (defaults $\alpha = 0.6$,
$\lambda = 0.5$, 5 rounds). Under positive damping the iteration is
a strict contraction on the probability simplex, so by Banach's
fixed-point theorem HBP converges geometrically to a unique
self-consistent belief (Theorem~2, Supplementary Note~S1). At
default hyperparameters the algorithm reaches four-decimal
convergence in $\approx$0.08~s per locus on 1000 Genomes chromosome~22,
approximately 20--30$\times$ faster than SuSiE and FINEMAP on the
same data.

\begin{figure}[htbp]
\centering
\includegraphics[width=0.95\textwidth]{fig2_hbp_gafm_schematic.pdf}
\caption{\textbf{HBP and GAFM, the two graph-native fine-mapping methods
benchmarked in this paper.}
(\textbf{a}) HBP performs message passing on a typed factor graph
with variant nodes (bottom), gene nodes (middle), and pathway nodes
(top). The upward pass aggregates variant beliefs into gene scores
weighted by eQTL strength, diffuses them once along the STRING
protein--protein interaction graph at a fixed 0.3 coupling, and
aggregates to the pathway layer. The downward pass reverses the
direction and produces a variant prior $\pi^{(t)}$. Each iteration
combines the prior with LD-deconvolved statistical evidence as
$b^{(t+1)} = \alpha \cdot \text{softmax}(u) + (1-\alpha) \pi^{(t)}$
and damps by $\lambda$ before renormalisation. Defaults are
$\alpha = 0.6$, $\lambda = 0.5$, $T = 5$ rounds. Under positive
damping the update is a strict $\ell_1$ contraction on the
probability simplex (Theorem~2, Supplementary Note~S1), guaranteeing
geometric convergence to a unique fixed point.
(\textbf{b}) GAFM combines LD-deconvolved unique-statistics
$u_i = z_i - |N_i|^{-1} \sum_{j \in N_i} R_{ij}^2 z_j$ with a
graph-derived functional score $z^{\mathrm{func}}_i$ (eQTL,
PPI-coupled gene activity, regulatory-element flag) via an
empirical-Bayes-tuned blend coefficient $\alpha$. The combined
score $s_i = \alpha u_i + (1-\alpha) z^{\mathrm{func}}_i$ is
softmaxed to per-variant PIPs, and (optionally) reweighted by a
4-component Wakefield mixture posterior to produce GAFM-MX,
HBP-MX and ENS. Under mild LD-decay assumptions the LD-deconvolved
unique-statistic $u_i$ is provably maximised at the causal variant
(Theorem~3). See Methods (``\nameref{methods:hbp}'',
``\nameref{methods:gafm}'') for the full algorithms.}
\label{fig:hbp_gafm}
\end{figure}

\subsection*{Graph-augmented fine-mapping (GAFM) wins weak signal 27--2 against SuSiE}

The regime that matters most for post-GWAS discovery is the one
where statistical signal alone is insufficient to resolve the causal
variant: sub-genome-wide-significant loci, rare-variant effects, and
complex LD blocks. We designed GAFM for this
regime. GAFM first LD-deconvolves the z-statistics
($u_i = z_i - \frac{1}{|N(i)|} \sum_{j \in N(i)} r^2_{ij} z_j$ with
$N(i) = \{j : r^2_{ij} > 0.3\}$), then combines the deconvolved
evidence with a multi-layer graph functional score $z_{\text{func}}$
via $s_i = \alpha u_i + (1-\alpha) z_{\text{func},i}$ with
$\alpha$ chosen by empirical Bayes on a held-out calibration set.
Under mild LD-decay assumptions the causal variant has the highest
unique-signal contribution (Theorem~3).

On 79 independent simulated loci with tissue-specific eQTL causal
variants and weak effect size ($\beta = 0.15$, $h^2 = 0.01$), GAFM
resolves 61/79 loci at rank~\#1 versus SuSiE's 45/79
(Figure~\ref{fig:weak_signal}). Mean causal-variant rank is 1.65
for GAFM versus 2.84 for SuSiE. Head-to-head on matched replicates, GAFM
beats SuSiE on 27, ties on 50, and loses on 2 (win ratio 13.5\,:\,1,
sign-test $p < 10^{-5}$). The runtime advantage persists: 0.07~s per
locus for GAFM versus 1.8~s for SuSiE.

The 27--2 headline is regime-specific, not universal. A 100-replicate
replication on 1000 Genomes chromosome~22 without tissue-specific
eQTL causal enrichment---$\beta = 0.2$, $h^2 = 0.02$, 30 rotating
centres across 17--46~Mb---converges to parity between GAFM (54\%
rank-\#1) and SuSiE (57\%), with GAFM retaining a $\sim$25$\times$
runtime advantage. The two results together bound the claim:
\textbf{GAFM's decisive statistical advantage emerges when annotation
informativeness is high} (tissue-specific eQTLs, moderate-impact
functional consequences, pathway co-enrichment); \textbf{under
uniform annotation conditions the advantage narrows to parity on
rank and persists only on speed.} This is the pattern that a
relational prior should produce, not a universally stronger
statistical test.

\begin{figure}[htbp]
\centering
\includegraphics[width=0.95\textwidth]{fig6_weak_signal_headline.pdf}
\caption{\textbf{GAFM beats SuSiE 27--2 head-to-head at weak signal with
tissue-specific eQTL priors.} 79 independent simulated loci on
1000 Genomes chromosome~22, with causal variants drawn from GTEx~v8
significant eQTLs at tissue-specificity $-\log_{10} p > 15$ in at
most two tissues; effect size $\beta = 0.15$, $h^2 = 0.01$.
(\textbf{a})~Distribution of causal-variant rank (lower is better):
GAFM mean 1.65, SuSiE mean 2.84. (\textbf{b})~Head-to-head scoreboard:
GAFM wins 27, ties 50, SuSiE wins 2; win ratio 13.5\,:\,1
(sign-test $p < 10^{-5}$). (\textbf{c})~Rank-\#1 rate: GAFM~77\%
(61/79) vs SuSiE~57\% (45/79). (\textbf{d})~Per-replicate scatter
of GAFM rank against SuSiE rank. Points above the diagonal are
SuSiE-favoured replicates (2/79); points below are GAFM-favoured
(27/79). Runtime 0.07~s for GAFM vs 1.8~s for SuSiE. At higher
annotation uniformity (see next subsection) the GAFM advantage
narrows to parity on rank.}
\label{fig:weak_signal}
\end{figure}

\subsection*{Graph and flat priors converge on accuracy; the graph wins on speed and on tissue-specific signal}

A common GraphGWAS interface integrates the six methods benchmarked
by Wu~\emph{et~al.}~2026\cite{wu2026genome}: SuSiE\cite{wang2020simple},
FINEMAP\cite{benner2016finemap}, SuSiE-inf\cite{cui2024improving},
FINEMAP-inf\cite{cui2024improving}, a Polyfun-proxy\cite{weissbrod2020polyfun}
that feeds the same eQTL priors to SuSiE via \texttt{prior\_weights},
and SBayesRC\cite{zheng2024leveraging}. We ran a uniform 6-method
head-to-head on 90 simulated 1000~Genomes chromosome~22 loci ---
30 strong-signal replicates ($\beta{=}0.5$, $h^2{=}0.10$, random
causal), 30 weak-signal ($\beta{=}0.2$, $h^2{=}0.02$, random) and 30
functional-causal ($\beta{=}0.2$, $h^2{=}0.02$, tissue-specific eQTL
restricted) --- with all six methods evaluated on the same loci with
the same simulation seeds (Figure~\ref{fig:baselines};
Supplementary Table~S2). At strong signal SuSiE-inf and FINEMAP-inf
lead marginally (23/30 rank-\#1, 77\%) over GAFM (22/30, 73\%) and
HBP (21/30, 70\%), with SuSiE/FINEMAP at 70--73\% --- consistent
with the small inf-extension advantage reported by Wu~\emph{et~al.}
At weak signal all six converge to 57--60\% rank-\#1 within three
percentage points; at functional-causal, all six converge to 37--40\%.
The graph methods complete each locus in 0.07--0.12~s versus
0.22--0.89~s for the inf-extension methods and 1.7--3.1~s for
vanilla SuSiE/FINEMAP --- a 5--40$\times$ speed advantage that holds
across all three regimes. When the eQTL prior is also supplied to
SuSiE via Polyfun-proxy on a parallel weak-signal benchmark, SuSiE's
rank-\#1 count rises from 10/20 to 11/20 and GAFM matches at 26$\times$
the speed. Combined with the 27--2 weak-signal result, the
head-to-head landscape resolves: \textbf{graph-native and flat-vector
priors converge on accuracy when the prior information is identical;
the graph wins decisively on speed, and on accuracy in regimes where
tissue-specific eQTL structure is informative.}

The 5--40$\times$ speed advantage is partly statistical (fewer
latent effects to fit) and partly a property of the data model.
Reconstructing each locus's multi-omics neighbourhood from
separate files (BED tracks for regulatory elements, GTF for genes,
TSV matrices for tissue-specific eQTLs) is index-bound;
the same information stored as a typed graph collapses the join
into a single index-free traversal that returns in milliseconds.
SBayesRC occupies a complementary problem (genome-wide multi-component
Bayesian mixture on $\sim$1.2~M HapMap~3 SNPs requiring $N \gtrsim
10^5$); HBP and GAFM target single-locus resolution at leads surfaced
by genome-wide methods.

\begin{figure}[htbp]
\centering
\includegraphics[width=0.95\textwidth]{fig3_hbp_vs_susie_finemap.pdf}
\caption{\textbf{HBP and GAFM match the rank-\#1 rate of five
Bayesian baselines at 5--40$\times$ the speed across all three
simulation regimes.}
Uniform 6-method head-to-head benchmark on 90 simulated loci from
1000 Genomes chromosome~22, 30 replicates per scenario, all six
methods evaluated on the same loci with the same seeds; per-method
numerics in Supplementary Table~S2. (\textbf{a})~Rank-\#1 rate of
the causal variant across the three scenarios (strong: $\beta{=}0.5$,
$h^2{=}0.10$, random causal; weak: $\beta{=}0.2$, $h^2{=}0.02$,
random causal; functional: $\beta{=}0.2$, $h^2{=}0.02$,
tissue-specific eQTL causal). At strong signal SuSiE-inf and
FINEMAP-inf lead marginally (23/30, 77\%) over GAFM (22/30, 73\%),
HBP (21/30, 70\%) and SuSiE/FINEMAP (21--22/30, 70--73\%). All six
methods converge to within three percentage points at weak signal
(57--60\%) and at functional (37--40\%).
(\textbf{b})~Mean causal-variant rank (log scale, lower better)
across the same 6$\times$3 matrix; the picture mirrors panel a but
makes the magnitude of mis-rankings visible (functional regime
ranks span 11--19 across methods).
(\textbf{c})~Per-locus runtime distributions (log scale) pooled
across all 90 simulations: HBP~0.08--0.12~s, GAFM~0.07--0.11~s,
SuSiE-inf~0.22--0.41~s, FINEMAP-inf~0.55--0.89~s,
SuSiE~1.7--2.2~s, FINEMAP~2.2--3.1~s. The graph methods are
5--40$\times$ faster than the inf-extension baselines and
20--30$\times$ faster than vanilla SuSiE/FINEMAP.
(\textbf{d})~Polyfun-proxy weak-signal benchmark ($h^2 = 0.02$, 20
replicates, separate experiment): when the eQTL prior used by
HBP/GAFM is also supplied to SuSiE via the \texttt{prior\_weights}
mechanism, SuSiE's rank-\#1 count rises from 10/20 (vanilla) to
11/20 (with prior) and GAFM matches at 11/20 with $\approx$26$\times$
less wall-clock time. SBayesRC, the sixth Wu~\emph{et~al.}\
baseline, addresses a genome-wide multi-component mixture problem at
$N \gtrsim 10^5$ and is reported separately in Supplementary
Table~S2 rather than co-plotted here.}
\label{fig:baselines}
\end{figure}

\subsection*{PIPs are calibrated, null false-positive rate is zero, and quality scales monotonically with $N$}

PIPs must mean what they say and must not fire under the null. On
100 null simulations (no causal variant, Gaussian phenotype) over
random 50-kb yeast loci, every method (HBP, GAFM, SuSiE, FINEMAP)
reports $\max$\,PIP~$<$~0.5 in every replicate, with mean
$\max$\,PIP at 0.003--0.029 --- matching to within 2\% the
softmax bound $\mathbb{E}[\max_i \pi_i] \leq
O(\sqrt{\log n}/n)$ proved in Supplementary Note~S1. On 200
calibration simulations across $h^2 \in \{0.02, 0.05, 0.10, 0.20\}$,
true-discovery rates track the reported PIP diagonal for SuSiE,
FINEMAP and GAFM; HBP is mildly conservative at mid-range PIPs but
when it assigns PIP $>$ 0.5 the variant is causal 99\% of the
time (Supplementary Figure~S4).

HBP fine-mapping quality scales monotonically with sample size:
on 1000 Genomes chromosome~22 subsampled at
$N \in \{500, 1{,}000, 2{,}000, 3{,}000\}$ ($h^2 = 0.10$, causal
minor allele frequency (MAF) 10--40\%, 30 replicates per $N$),
mean causal PIP rises
0.54~$\rightarrow$~0.58~$\rightarrow$~0.59~$\rightarrow$~0.77, with
73\% of replicates at rank~1 by $N = 3{,}000$ (Supplementary
Figure~S5). The linear trend extrapolates beyond the training
regime, supporting the biobank-scale claim tested directly on
Pan-UK Biobank in the next subsection.

\subsection*{HBP fine-maps Pan-UK Biobank across four ancestries with no algorithmic change}\label{sec:cross_anc}

LD structure differs substantially across ancestries and many causal
variants are ancestry-private. On 1KG chromosome~22 restricted to
the three largest superpopulations (EUR $n=503$, AFR $n=661$, EAS
$n=504$, F1 simulations at $h^2 = 0.05$, 30 replicates each), HBP
rank-\#1 ordering is \textbf{AFR 53\% $>$ EAS 27\% $>$ EUR 17\%}
(mean PIPs 0.35, 0.21, 0.08; Figure~\ref{fig:cross_anc}a) --- the
expected direction, since AFR's older, less-correlated LD produces
fewer indistinguishable proxies per locus. HBP's precision holds
across all three ancestries.

A sumstats-only entry path consumes Pan-UK
Biobank\cite{karczewski2024pan,bycroft2018uk} summary statistics
streamed via \texttt{tabix}\cite{li2011tabix} over HTTPS, with
GRCh37$\to$GRCh38 liftover to align with the 1KG reference panel.
Applied to four canonical multi-ancestry loci ---
\emph{FTO}/BMI\cite{frayling2007common},
\emph{APOA5}/triglycerides\cite{do2013common},
\emph{LDLR}/LDL-C\cite{teslovich2010biological},
\emph{HMGA2}/height\cite{weedon2008genome,yengo2022saturated} ---
across four Pan-UKB ancestries (EUR $N=420{,}531$; CSA $8{,}876$;
AFR $6{,}636$; EAS $2{,}709$), three of the four EUR cells (BMI,
LDL, height) resolve to a \textbf{single-variant credible set at
PIP $=$ 1.000}. \emph{HMGA2}/height is cross-ancestry robust: CSA
PIP $=$ 1.000 at $N=8{,}876$ and EAS PIP $=$ 0.97 at $N=2{,}709$ ---
a 145$\times$-smaller sample than EUR. \emph{APOA5}/triglycerides
shows two-ancestry convergence on the same indel
(EUR PIP $=$ 0.51, CSA PIP $=$ 0.73). The same codebase and the
same algorithm; no EUR-specific tuning is needed.

\begin{figure}[htbp]
\centering
\includegraphics[width=0.95\textwidth]{fig9_cross_ancestry.pdf}
\caption{\textbf{HBP fine-mapping holds across four Pan-UKB
ancestries and matches the simulated 1KG cross-ancestry trend on
real biobank data.}
(\textbf{a})~Controlled simulation: 1KG chromosome~22 restricted to
three superpopulations (EUR $n{=}503$, AFR $n{=}661$, EAS $n{=}504$);
F1 simulations at $h^2 = 0.05$, 30 replicates per ancestry. HBP
rank-\#1 rate AFR~53\% $>$ EAS~27\% $>$ EUR~17\% reflects AFR's older,
less-correlated LD. (\textbf{b})~Real-biobank demonstration: HBP
sumstats-only fine-mapping on four canonical loci (FTO~/~BMI,
APOA5~/~triglycerides, LDLR~/~LDL-C, HMGA2~/~height) across four
Pan-UKB ancestries (EUR $N{=}420{,}531$; CSA $8{,}876$; AFR $6{,}636$;
EAS $2{,}709$). Three EUR cells resolve to a single-variant 95\%
credible set at PIP~=~1.0. Height/HMGA2 is cross-ancestry robust:
CSA PIP~=~1.0, EAS PIP~=~0.97 at 145$\times$ smaller $N$ than EUR.
Triglycerides/APOA5 shows two-ancestry convergence on the same indel.
No EUR-specific tuning; no algorithmic change between the simulated
and real-biobank panels.}
\label{fig:cross_anc}
\end{figure}

\subsection*{The same pipeline fine-maps yeast, \emph{Arabidopsis}, rice and human at biobank and germplasm-bank scale}
\label{sec:multispecies}

Fine-mapping causal variants is hard for distinct reasons across the
scales human medical genetics, model-organism crosses, and crop
germplasm operate at: in human cohorts, signal-to-noise is high but
LD is dense and reference panels are ancestry-specific; in
inbred-yeast and \emph{Arabidopsis} panels, LD is exceptionally long
and pathway/PPI annotation is sparse; in crop germplasm such as the
3{,}000 Rice Genomes panel, accessions span large between-population
$F_{ST}$ and most traits have only ordinal phenotypic scoring with
no curated causal-gene catalogue. The same fine-mapping pipeline,
unchanged except for the loaded annotation graph, was applied to all
three regimes (Supplementary Tables~S5--S6).

\paragraph{Rice (3{,}000 Rice Genomes, 18 IRRI ordinal traits + 4 grain weight/shape traits).}
The rice 3kRG analysis comprises two complementary GWAS passes on
the same 3{,}024-accession genotype substrate. The breadth pass is
a whole-genome GWAS on 18 IRRI Standard Evaluation System ordinal
trait codes\cite{wang2018genomic} ($\lambda_{GC}$: 0.50--2.07,
median 1.09; per-trait detail in Supplementary Table~S5) used as
the substrate for fine-mapping the top three genome-wide-significant
plus one suggestive independent lead per trait (72 total) within
$\pm$250~kb windows against a rice multi-omics graph (RAP-DB
variant--gene + Ren 2023 gene--pathway curation\cite{ren2023rice}
+ RicePPINet gene--gene\cite{liu2017riceppinet}). The depth pass
is a 4-trait GWAS on continuous grain weight (TGW) and grain shape
(GL, GW, RLW = GL/GW) phenotypes from Niu \emph{et al.}\ 2021
\cite{niu2021grain} on the same panel ($N{=}1{,}847$ for TGW;
$N{=}2{,}453$ for GL/GW/RLW; $h^2$ 0.88--0.93 in Niu's
variance-component decomposition; $\lambda_{GC}$: 1.41--1.78 under
PC1--PC10 correction, reflecting the same XI/GJ population
structure as the IRRI scan; pairwise Pearson correlations match
Niu's Fig.~1b to within $\Delta r{<}0.04$). The grain pass provides
the strongest crop-genomics validation in the manuscript: Niu
\emph{et al.}\ 2021 reports 21 stable QTNs identified across 3
years of trials by compressed-MLM analysis on the same panel,
giving a literature-curated, panel-matched ground truth not
available for the 18-trait IRRI scan.

GAFM delivers a single-variant 95\% credible set with PIP~$\geq$~0.99
on 11 of 72 leads (15\%). \textbf{Twenty-four of 72 leads (33\%)
fine-map to within 250~kb of a Ren-2023 catalogued causal gene}, with
14 of 18 traits contributing at least one validated locus. Notable
hits include the rice red-pericarp gene \emph{Rc} only \textbf{4.5~kb}
from the spikelet-pericarp-colour lead at Chr7:6.07~Mb;
\textbf{\emph{TGW6}} 12.3~kb from the leaf-length-type Chr6:25.11~Mb
lead; the awn-development gene \emph{GAD1/RAE2} 4.4~kb from the
awning-presence Chr8:23.98~Mb lead ($p=4.4\times10^{-90}$,
PIP~$=$~0.999); grain-shape genes \textbf{\emph{An-1}}
($240$~kb, blade-colour Chr4:16.47~Mb, PIP~$=$~1.00) and
\textbf{\emph{BG1}} ($68.7$~kb, spikelet-fertility Chr3:4.10~Mb);
the panicle-development gene \emph{OsER1} ($62$~kb,
apiculus-colour Chr6:5.31~Mb, PIP~$=$~0.77); and the seed-dormancy
regulator \emph{OsVP1} ($176$~kb, panicle-type Chr1:39.54~Mb).
The two highest-confidence pigment-related leads validate against
the same axis of variation captured by the IRRI ordinal scoring
(spikelet pericarp at Chr7:6.07~Mb / \emph{Rc} and apiculus colour
at Chr10:13.32~Mb / \emph{BRD2/LTBSG1}, $4.5$~kb and $28.5$~kb
respectively, both at PIP $\geq 0.99$). The cross-species
fine-mapping claim for rice is restricted to the 14 / 18 IRRI
traits with $\lambda_{GC} \in [0.85, 1.20]$ (well-calibrated
subset); the remaining four traits with $\lambda_{GC} \in
[1.6, 2.1]$ show residual subpopulation
confounding ($F_{ST} \approx 0.5$ between XI and GJ subpopulations
under 10-PC correction); the eight doubly-resolved loci with
HBP CS $=$ GAFM CS $=$ 1 are all in well-calibrated traits
($\lambda_{GC} \in [0.50, 1.38]$).

\textbf{Grain weight + shape pass.}
The Niu \emph{et al.}\ 2021 phenotype panel ($N{=}2{,}453$ for grain
shape; $N{=}1{,}847$ for grain weight) yields 4{,}454 independent
genome-wide-significant leads across the four traits at 250\,kb
clumping. \textbf{Twenty of 21 stable QTNs identified by Niu
\emph{et al.}\ 2021 are recovered within $\pm$100\,kb of an
independent lead at $p \leq 5\times10^{-8}$} (95\% catalogue
recovery at GWAS lead-call resolution; the only miss is qGL10 at
Chr10:13.62\,Mb, originally reported by Niu at
$p{=}1.19\times10^{-7}$ in compressed-MLM, which falls below our
$p{<}10^{-5}$ suggestive cutoff under the linear-model + 10-PC
pipeline). Eleven of the recovered QTNs match the lead position
to within $\leq$10\,kb, including the \emph{GS3}
(Chr3:16{,}733{,}441) and \emph{qSW5/GW5} (Chr5:5{,}371{,}529)
hubs that drive grain length, width and ratio simultaneously.

\textbf{Nine-method fine-mapping at 41 leads.} We fine-mapped a
literature-anchored set of 41 leads (top 5 GW-significant
$+$ top 2 suggestive per trait, augmented by every lead falling
within $\pm$100\,kb of a Niu QTN; $\pm$100\,kb LD windows;
MAF~$\geq$~0.05). All nine methods receive
$\lambda_{GC}$-deflated z-scores (the standard genomic-control
correction; $z \to z/\sqrt{\lambda_{GC}}$). The panel comprises
the original GAFM and HBP; their mixture-prior posterior
reweightings on the LD-deconvolved deflated z-scores (GAFM-MX,
HBP-MX); the
mean-of-PIPs ensemble of the two augmented methods (ENS); and
the four canonical baselines SuSiE \cite{wang2020simple},
SuSiE-inf, FINEMAP-inf \cite{cui2024improving}, and SBayesRC
\cite{zheng2024leveraging}. The SBayesRC LD reference is a
custom 23-block 3kRG-specific eigen-decomposed reference covering
all 41 leads (188\,K MAF~$\geq$~0.05 variants; gctb 2.05beta +
R \texttt{LDstep1}--\texttt{LDstep4}; ${\sim}$10\,min compute).

At the strict variant-level criterion (the QTN's chr:pos in the
95\% credible set within $\pm$10\,kb), per-method recovery of the
21 Niu QTNs is \textbf{SBayesRC 20/21 (95\%, 16 exact-position
matches)} $>$ \textbf{GAFM-MX = HBP-MX = ENS = SuSiE 12/21 (57\%;
GAFM-MX/HBP-MX/ENS 7 exact, SuSiE 5 exact)} $>$ HBP 10/21 (48\%)
$>$ FINEMAP-inf 8/21 (38\%) = SuSiE-inf 8/21 = GAFM 8/21 (38\%).
The mixture prior posterior reweighting closes a substantial
fraction of the gap to SBayesRC: GAFM-MX gains 19 percentage
points over GAFM (38\%$\to$57\%) and HBP-MX gains 9 percentage
points over HBP (48\%$\to$57\%). At the \textbf{stricter
top-1-PIP exact-position metric (the QTN is the single
highest-PIP variant in the credible set), GAFM-MX, HBP-MX, and
ENS each achieve 10/21 = 47.6\% --- the highest of any method,
exceeding SuSiE's 6/21 (28.6\%) and SBayesRC's 3/21 (14.3\%)}. For the 7 NEW candidate genes (locus-level), GAFM-MX,
HBP-MX, ENS, SuSiE, SuSiE-inf and SBayesRC all reach
6/7 = 85.7\%; HBP and FINEMAP-inf reach 4--5/7 (57--71\%); GAFM
3/7 (43\%). Ren 2023 catalogue gene recovery sits at 16--18 of
the $\sim$118 grain-shape genes across all nine methods.

The reframing shows that the canonical fine-mapping baselines
(SuSiE family) and the graph-augmented GAFM/HBP can be brought
to within striking distance of SBayesRC on the strict variant-level
metric by importing a single ingredient that SBayesRC bakes into
its MCMC --- the 4-component spike-and-slab effect-size prior ---
as a post-hoc Bayes-factor reweighting step. The remaining gap
to SBayesRC's 95\% recovery (28 percentage points) reflects
SBayesRC's full eigen-decomposed LD model and its joint-MCMC
inference across all 188\,K variants in the 23-block window;
GAFM-MX and HBP-MX retain the 200--700$\times$ speed advantage
of the original GAFM/HBP (median 0.027 s and 0.029 s per locus,
respectively, including the mixture-prior reweight; vs.\ 19.2 s
SuSiE, 5.4 s SuSiE-inf, 4.7 s FINEMAP-inf, ${\sim}$10\,min
SBayesRC genome-wide). Per-method, per-tier recovery is in
Supplementary Table~S8; per-locus credible sets in Supplementary
Table~S9.

The depth pass thus extends the cross-species claim from
``IRRI-ordinal traits with $\lambda_{GC} \in [0.85, 1.20]$''
to ``four continuous grain weight/shape traits whose
literature-curated 21-QTN catalogue we recover at 95\% rate at
GWAS-lead resolution and at $\sim$48\% rate at variant-level
resolution against a panel-matched ground truth without re-tuning
the pipeline'' --- the strongest crop-genomics benchmark in the
manuscript. The unusually inflated $\lambda_{GC}$ values for this
pass (1.41--1.78) are consistent with the well-documented XI/GJ
subpopulation structure of the 3kRG panel under PC1--PC10
correction; SuSiE and SuSiE-inf give CS$=$1 on every locus
(40--41 / 41), which is consistent with the inflation-induced
single-strongest-variant collapse and is the same caveat raised
for the IRRI 18-trait pass at $\lambda_{GC} \in [1.6, 2.1]$
(\S\ref{sec:multispecies} above). GAFM/HBP are more conservative
under inflation (median CS$=$1{,}823 and 814 respectively at the
secondary leads); the variant-level recovery rate is therefore
the most defensible cross-method comparison metric here.

\paragraph{Yeast (1011 Genome panel, 35 growth traits).}
Across 35 grammar-corrected growth-trait sumstats\cite{peter2018genome,
bloom2015genetic} fine-mapped against an SGD-derived
gene--GO-slim graph (1.14~M variants, 60\% of the panel), 245 leads
were fine-mapped at $\pm$15~kb windows. Yeast LD is strong enough that
no locus reaches GAFM CS $=$ 1, but \textbf{HBP narrows the credible
set on 244 of 245 leads (100\%)}, with 5 loci reaching PIP $\geq 0.5$.
Forty-four of the 245 leads (18\%) lie within 250~kb of a
literature-curated trait-relevant gene (e.g.\ TUP1/SPT7 for ethanol,
CUP1/CUP2 for copper resistance, the GAL pathway for galactose,
ENA1/HOG1 for salt response, ERG/PDR for fluconazole).

\paragraph{Arabidopsis (1001 Genomes, 14 priority traits).}
A TAIR10 + Plant Reactome + STRING-\emph{Ath} graph (covering
1.06--3.99~M variants across the five chromosomes; 20{,}192 genes
with high-confidence PPI partners after UniProt-orthologue
resolution) supports a multi-trait \texttt{plink2} GWAS on 14
priority phenotypes (11 AraGWAS-Bonferroni-validated traits +
flowering-time traits FT10, FT16 and FLC). Of the 54 fine-mapped
leads, three reach GAFM CS $=$ 1 with PIP $\geq 0.98$ (a seed-storage
trait at three independent chr4 and chr5 loci); HBP narrows the
credible set on 29 of 54 (54\%). Validation against AraGWAS plus
canonical flowering-time genes (FRI, FLC, CONSTANS, GIGANTEA, VIN3,
FT, DOG1, HKT1) gives 3 of 54 (6\%); the apparently low rate
reflects the chr-4-only coverage of the AraGWAS Bonferroni file
rather than poor fine-mapping.

\paragraph{Human (Pan-UKB EUR, four phenotypes, all 22 autosomes).}
A genome-wide GTEx + STRING graph (\textbf{4.59~M annotated
variants}; 49 GTEx tissues + STRING v12 PPI at score $\geq$ 700;
GENCODE v47 gene model) supports tabix-over-HTTPS Pan-UKB summary
statistics with on-the-fly GRCh37$\to$GRCh38 liftover and per-chromosome
1KG NYGC LD references (3{,}202 samples, GRCh38). \textbf{Three
hundred and twenty-one leads were fine-mapped across BMI, LDL
direct, and triglycerides}; the Pan-UKB Height sumstats are heavily
flagged \texttt{low\_confidence} genome-wide under the
500~kb-clustering filter we applied and returned no leads per
chromosome (the locus-targeted Height/HMGA2 query reported in
\S``\nameref{sec:cross_anc}'' bypasses this clustering by querying a specific
500~kb window directly).

\textbf{Genome-wide GAFM recovered the textbook canonical
lipid-metabolism genes at single-variant resolution, without any
annotation prior:} \textbf{\emph{LDLR}} (LDL chr19:11.08~Mb, GAFM
CS $=$ 1, PIP $=$ 1.00, 12.7~kb from the Brown \& Goldstein 1986
prototype gene); \textbf{\emph{APOE}} (LDL and triglycerides
chr19:44.81/44.91~Mb, PIP $=$ 0.60--0.77);
\textbf{\emph{LPL}} (chr8:20.0~Mb, GAFM CS $=$ 1, PIP $=$ 0.9998);
\textbf{\emph{GCKR}} (chr2:27.5~Mb, GAFM CS $=$ 2, PIP $=$ 0.92,
\textbf{0~bp} from the gene body); \textbf{\emph{ANGPTL3}}
(chr1:62.97~Mb, PIP $=$ 0.81, \textbf{0~bp});
\textbf{\emph{ANGPTL4}} (chr19:8.4~Mb); and the
\textbf{\emph{APOA5} cluster} (chr11:117.0~Mb, PIP $=$ 0.65). Eight of
321 leads reach GAFM CS $=$ 1; 47 reach PIP $\geq 0.5$;
the strong-PIP subset (GAFM CS $\leq 5$ \emph{and} PIP $\geq 0.5$)
matches a curated chr-22 / GWAS-Catalog literature catalogue at 9 of
10 leads. Across the genome-wide scan, HBP and GAFM produce identical
credible sets on every one of 321 leads --- a finding to which we
return below.

\subsection*{A non-uniform per-variant prior, not uniform high coverage, lets the graph break LD ties}
\label{sec:heterogeneity}

The cross-species data above show 287 of 692 leads (41\%)
fine-mapped tighter under HBP than under flat-prior GAFM, but the
distribution of that gain across species is striking and
non-uniform: yeast 100\%, Arabidopsis 54\%, rice 19\%, human 0\%
genome-wide. The naive interpretation that "more annotation density
helps fine-mapping" is therefore \emph{wrong} as a general rule:
the human cache is by far the densest of the four (4.59~M
annotated variants, 49 tissues, near-saturating GTEx + STRING
coverage) yet shows zero HBP narrowing relative to GAFM. The
explanation is that the human cache is also nearly \emph{uniform}
per variant --- almost every annotated variant has the same
neighbourhood structure (one eGene plus tens of STRING partners).
Random graph priors with low per-variant variation cannot break
LD ties that flat-prior fine-mapping has not already broken.

\paragraph{Ablation experiment.} We tested this directly on the
3{,}000 Rice Genomes panel by randomly subsampling the rice
multi-omics cache at five retention fractions (10\%, 25\%, 50\%, 75\%,
100\%) $\times$ 10 bootstrap seeds, and re-running HBP on each of
five simulated grain-quality causal loci with curated ground truth
(Figure~\ref{fig:ablation}). Three regimes emerged: at
already-resolvable loci (baseline PIP $\geq 0.01$), more coverage
has no visible effect (the causal stays at rank 2);
\textbf{at borderline loci (baseline PIP $\sim 10^{-3}$), more
coverage is dramatically helpful} --- BG1's causal-variant rank
improves from 154 at 10\% coverage to 4 at 100\% (a $38.6\times$
gain) and TGW6's improves from 65 to 3 ($21.6\times$). At
underpowered loci (baseline PIP $\leq 10^{-4}$), more coverage
\emph{degrades} the causal's rank, because added annotation mass is
dispersed across non-causal variants. \emph{Heterogeneity}, not
density, drove the response: rice's hand-curated 269-gene pathway
labels and RicePPINet edges produced cross-seed variance in HBP
output (rank standard deviation 7--18 at borderline loci),
whereas the uniform human cache produced exactly zero variance
across seeds.

\paragraph{Direct intervention.} If the heterogeneity-not-density
principle is causal, then injecting per-variant heterogeneous
prior information into the cache --- without changing GWAS
sumstats, LD, or the fine-mapping algorithm --- should turn the
human cache from inert (HBP $=$ GAFM) into informative (HBP narrows
beyond GAFM). We extended the cache schema with an optional
per-variant continuous score that HBP and GAFM mix into the
final posterior (Methods). For human, the smallest available
heterogeneous layer was an ENCODE cCRE class flag\cite{encode2020expanded}
(926{,}535 elements across nine classes:
proximal/distal enhancer-like, promoter-like, CTCF-bound, etc.;
$\approx$10--15\% of GTEx-annotated variants gain a class-weighted
score). For Arabidopsis we used TAIR10 CDS-vs-gene-body class;
for rice, snpEff impact class (HIGH/MODERATE/LOW/MODIFIER); for
yeast, BIOGRID physical interactions plus ORF feature class.

\textbf{The intervention transforms the human result and confirms
the principle in three of four species} (Table~\ref{tab:intervention}).
On the same 321 human Pan-UKB leads, adding the cCRE class flag
alone takes HBP narrower than GAFM on \textbf{282 of 321 leads (88\%)},
up from 0 of 321 (0\%) without it --- a 88-percentage-point
intervention effect on a fixed locus set. The same direction holds
for Arabidopsis (54\%$\to$94\%) and rice (19\%$\to$78\%); yeast was
already saturated at 100\%. Across all 692 leads, HBP narrows
beyond GAFM on 287 (41\%) before the intervention and 634 (92\%)
after.

\begin{figure}[htbp]
\centering
\includegraphics[width=0.85\textwidth]{fig_ablation_rice.pdf}
\caption{\textbf{Multi-omics coverage helps borderline-resolvable
loci most, has no effect on already-resolvable loci, and slightly
degrades unresolvable ones.}
For each of five curated grain-quality causal loci, HBP
was re-run with 10--100\% of multi-omics edges randomly retained
(10 bootstrap seeds per non-100\% fraction). The rank of the true
causal variant (log scale, lower is better) is plotted against
retained coverage. Three regimes emerge: \emph{resolvable} loci
(green, dashed) are already confident from sumstats and LD alone
and are coverage-insensitive; \emph{borderline} loci (blue, solid)
improve 20--40$\times$ in rank as coverage grows --- the regime
where the relational prior adds the most value;
\emph{underpowered} loci (red, dotted) cannot be rescued by
annotation and degrade slightly under dense priors.}
\label{fig:ablation}
\end{figure}

\begin{table}[htbp]
\centering
\footnotesize
\setlength{\tabcolsep}{4pt}
\caption{\textbf{Adding a non-uniform per-variant prior to the
multi-omics cache turns HBP from inert to informative across four
species.} ``HBP narrows tighter than GAFM'' is the percentage of
leads, per species, for which HBP returns a 95\% credible set with
strictly fewer variants than the flat-prior GAFM credible set on
the same locus data. ``Baseline cache'' = the multi-omics cache
without the per-variant continuous \texttt{prior\_score} field;
``+non-uniform prior'' = the same cache with a species-specific
per-variant continuous prior added (ENCODE cCRE 9-class flag for
human, TAIR10 feature class for \emph{Arabidopsis}, snpEff impact
class for rice, ORF-feature class plus BIOGRID for yeast). The
intervention does not alter GWAS sumstats, LD references, or the
fine-mapping algorithm. The human row is the cleanest demonstration:
a uniform GTEx + STRING cache leaves HBP exactly equal to GAFM on
every one of 321 leads, but adding a class-weighted ENCODE cCRE
flag flips that rate from 0\% to 88\% on the same lead set.}
\label{tab:intervention}
\begin{tabular}{@{}l@{\hspace{6pt}}r@{\hspace{6pt}}r@{\hspace{6pt}}r@{\hspace{6pt}}r@{}}
\toprule
\textbf{Species} & \textbf{Loci}
                 & \shortstack[r]{\textbf{Baseline cache}\\\textbf{HBP $<$ GAFM CS}}
                 & \shortstack[r]{\textbf{+non-uniform prior}\\\textbf{HBP $<$ GAFM CS}}
                 & $\boldsymbol{\Delta}$ \\
\midrule
Human       & 321 & 0 / 321 (0\%)   & 282 / 321 (88\%) & \textbf{$+$88\,pp} \\
Arabidopsis &  54 & 29 / 54 (54\%)  & 51 / 54 (94\%)   & $+$41\,pp \\
Rice        &  72 & 14 / 72 (19\%)  & 56 / 72 (78\%)   & $+$58\,pp \\
Yeast       & 245 & 244 / 245 (100\%) & 245 / 245 (100\%) & saturated \\
\midrule
\textbf{Total} & 692 & \textbf{287 (41\%)} & \textbf{634 (92\%)} & \textbf{$+$50\,pp} \\
\bottomrule
\end{tabular}
\end{table}

A pre-registered placebo-prior control on chr~22 (Sup Note~S2)
shows that on randomly-placed simulated causal variants, an
informative cCRE prior and a permuted-marginal placebo prior
produce \emph{identical} HBP-tighter rates (both 13\,\%, vs 10\,\%
baseline). The +3\,pp gain in simulation is therefore fully
captured by heterogeneity per se; the genome-wide
0\,\%\,$\to$\,88\,\% flip on 321 Pan-UKB lead loci (Table~1) must
depend on enrichment of real disease-associated leads in
regulatory regions rather than on the graph prior providing
variant-specific biological information beyond what any
non-uniform prior would produce. With this caveat, the
intervention results bound the claim to a
\emph{necessary-condition} statement:
a uniform high-coverage cache breaks no LD ties, while introducing
a per-variant non-uniform signal of biological structure (here, the
ENCODE 9-class regulatory-element flag) flips HBP from inert to
informative on 88\% of human leads. The current experiment
contrasts a uniform cache against an informative non-uniform
prior; a placebo-prior arm in which \texttt{prior\_score} is
permuted across variants (preserving the marginal distribution but
randomising which variants receive which class) is the immediate
next test, and is required before the claim that the gain
specifically reflects \emph{biological} discriminative information
rather than any heterogeneous prior is fully established.
Continuous-score replacements --- combined-annotation PHRED scores,
allele-specific transcription-factor-binding $\Delta$-PWM,
deep-learning chromatin-effect predictions, and species-specific
cistromes --- are a natural further direction.

